% Generated draft: requires mathematical and source review.
\documentclass[11pt]{article}
\usepackage[margin=25mm]{geometry}
\usepackage{fontspec}
\setmainfont{Latin Modern Roman}
\usepackage{amsmath,amssymb,mathtools,unicode-math}
\setmathfont{Latin Modern Math}
\usepackage{enumitem,needspace,xurl,hyperref}
\setlength{\parindent}{0pt}
\setlength{\parskip}{5pt}
\setlength{\emergencystretch}{3em}
\newcommand{\N}{\mathbb{N}}
\newcommand{\Z}{\mathbb{Z}}
\newcommand{\Q}{\mathbb{Q}}
\newcommand{\R}{\mathbb{R}}
\newcommand{\C}{\mathbb{C}}
\newcommand{\F}{\mathbb{F}}
\newcommand{\A}{\mathbb{A}}
\newcommand{\PP}{\mathbb{P}}
\newcommand{\E}{\mathbb{E}}
\newcommand{\eps}{\varepsilon}
\newcommand{\dd}{\,\mathrm{d}}
\begin{document}
\section*{UnsolvedMath ID: 800102}
\textbf{Source problem code:} catalogue entry
\subsection{Definitions and mathematical statement}
Write \(S^1=\{(x,y)\in\R^2:x^2+y^2=1\}\). The source renders the circle as \(\texttt{**S**}^1\); interpreting this as \(S^1\) is a typographical normalization requiring confirmation. An action of a group \(G\) on a set \(X\) assigns to each \(g\in G\) a bijection \(\rho(g):X\to X\), with \(\rho(gh)=\rho(g)\circ\rho(h)\) and \(\rho(e)=\operatorname{id}_X\). The supplied record does not specify \(G\), whether \(G\) carries a topology, or the continuity requirements for its actions.

For smooth one-dimensional foliations \(\mathcal F_1,\mathcal F_2\) of \(\R^2\), transversality means
\[
T_p\mathcal F_1\oplus T_p\mathcal F_2=T_p\R^2
\qquad\text{for every }p\in\R^2.
\]
An action preserves a foliation if every acting map sends each leaf onto a leaf of that foliation. Whether the source requires preservation of each foliation individually, and whether its foliations are smooth, remain unspecified.

\textbf{Existence question (supplied formulation).} Does there exist an action \(\varphi\) on \(S^1\) which is induced by a topological action \(\Phi\), but which is not induced by any smooth action \(\Psi\) on \(\R^2\) preserving two transverse foliations?

\textbf{Unresolved qualifications.} The word ``induced'' is not defined in the supplied record. In particular, no construction associates a circle action to a plane action, and no criterion specifies whether the resulting circle action must equal \(\varphi\) or merely be conjugate to it. The record also does not explicitly give the space on which \(\Phi\) acts or state whether \(\Phi\) must preserve two transverse foliations. Reading these as plane actions preserving such foliations is a possible interpretation, not an established hypothesis. It is unspecified whether the foliations are fixed in advance or may vary with the realizing action. The acting group, the quantification over that group, orientation conditions, and the intended smoothness requirements are likewise absent. Consequently, the displayed existence question preserves the source wording but does not yet determine a fully specified mathematical assertion.
\subsection{Short English statement}
The source asks whether a circle action can arise from a topological action while admitting no realization by a smooth plane action preserving two transverse foliations. The acting group and the construction that induces the circle action are missing, so the intended realization problem requires clarification before a precise formal statement can be fixed.
\subsection{Sources}
\begin{itemize}
\item Local UnsolvedMath registry, permanent ID \(800102\), source problem code \texttt{catalogue entry}, title ``Foliations and Diffeomorphism Groups --- Question 43.'' The supplied record lists Nariman and Yazdi as proposers and gives the year as 2025.
\item Nariman and Yazdi, \emph{Problem List on Foliations and Diffeomorphism Groups} (2025), Question 43. Supplied source URL: \url{https://arxiv.org/abs/2503.18274}. The provenance records access on 2026-10-16 and extraction by \texttt{arxiv\_latex}; these are supplied provenance claims, not independent source verification.
\item The registry labels the problem ``open'' and records creation and update timestamps of \texttt{2026-10-16T14:00:00Z}. This is supplied curation, not an independently verified assessment of mathematical status. Supplied upstream revision: \texttt{d371ee917ac9647fbb06cd6fc529deec6f3dda82}.
\end{itemize}
\end{document}
